\section{Threat model and goals}\label{sec:threat}

\paragraph{Parties.} The \emph{user} holds a passkey for the account and,
once it is lost, a new device. The \emph{mail provider} signs the user's
outgoing mail and publishes its keys in DNS. The \emph{DKIM key registry} is
a contract governed by a $k$-of-$n$ multisig and fed by independent DNS
attesters. The wallet's \emph{operator} serves the wallet, relays
transactions, controls the recovery address's domain, and runs the alert
service: an untrusted host and a \emph{sealed vault}, an attested enclave
that alone reads alert addresses and inbound mail. Alerts leave through a
third-party email-sending service. The account's other guardians are
whatever its recovery policy lists.

\paragraph{Adversaries.} We consider eight, alone or together. The
\emph{inbox thief} controls the victim's inbox, Sent folder included, and
wants the account. The \emph{social engineer} can get the user to send an
email with a chosen Subject to a chosen address and to read out or type the
recovery code, for example posing as support. The \emph{registry operator}
is registry members below the threshold, or a threshold of them adding a key
the provider never published. The \emph{relayer} or \emph{front-runner}
drops, delays, reorders or copies the user's transactions, or submits a copy
first, for another statement, attempt or account. The \emph{mail provider}
sees the user's mail and holds the signing key; it wants to learn which
account its user recovers. The \emph{network observer} sees the chain and
every pending transaction. The \emph{malicious frontend} is a misbehaving
wallet build: at recovery an attacker like any other client, at enrollment
trusted (A7). The \emph{operator's staff} read and change the alert host's
process, storage and logs, and the relayer, and want users' addresses; they
do not alter the recovery domain's mail routing (its MX, MTA-STS policy or
mail exchanger's TLS keys) or the wallet build users run (A10, A11), and
against staff who do, nothing built detects it (\S\ref{sec:limits}). Any adversary may open recovery
attempts against any account, send arbitrary email and run its own mail
domain. None breaks the cryptographic assumptions, reads the enclave's
memory, or holds two of the vault key's three shares.

\paragraph{Assumptions.} Table~\ref{tab:assumptions} lists what the
properties rest on; each is a ledger entry with status \claimtag{assumed}.
A1 is split: unforgeability (A1a) is cryptographic; sender authentication
(A1b) is a property of each provider's submission path, which
Appendix~\ref{app:providers} reviews from public sources, with every gap
open; a controlled submission test from real accounts is a check before
mainnet. If A1b fails at a provider, for example by checking only the first
of several \texttt{From} fields and signing the last, the factor there reduces to knowing
the victim's address and $\code$, which still open the commitment, under
the same delay, quorum and alerts\claim{email-factor-soundness}. A9
separates two clocks: the windows this design adds are measured on ledger
close time and keep their wall-clock length whatever the network's ledger
rate; Perch's delay and the statement's expiry are ledger counts.

\begin{table}[t]
  \caption{Assumptions. Each is \claimtag{assumed} in the claim ledger.}
  \label{tab:assumptions}
  \small
  \begin{tabularx}{\columnwidth}{@{}l>{\raggedright\arraybackslash}X@{}}
    \toprule
    Id & Assumption \\
    \midrule
    A1a & RSASSA-PKCS1-v1\_5 with SHA-256 is unforgeable for 2048-bit moduli,
          and providers keep their keys secret\claim{dkim-unforgeability}. \\
    A1b & A provider signs a \texttt{From} naming one of its mailboxes only
          for that mailbox's holder: it refuses two \texttt{From} fields or
          authenticates the one it signs, sending as another in-domain
          address needs proof of control, and a reassigned address is
          useless without $\code$\claim{provider-authenticates-sender}. \\
    A2, A3 & SHA-256 and Poseidon2 are collision resistant, and no known
          input maps to zero under Poseidon2; for hiding, Poseidon2 is a
          random function\claims{hash-collision-resistance,commitment-hiding,blinded-subject-hiding}. \\
    A4 & UltraHonk in Perch's zero-knowledge flavor is zero knowledge and
          weakly simulation extractable for adversaries that see honest
          proofs (stated, not established), and the pinned toolchain compiles
          and proves the circuit as
          written\claims{proof-system-security,toolchain-correctness}. \\
    A5 & Every key the registry lists for a domain was published by that
          domain's provider\claim{registry-correctness}. \\
    A6 & $\code$ and $\blind$ come from a secure generator and stay with the
          user and the browser\claim{secret-generation}. \\
    A7 & The enrolling wallet commits to the real enrollment email's address
          and pins only keys it checked\claim{honest-enrollment}. \\
    A8 & The enclave keeps its memory from its host, and at most one of the
          three key custodians is malicious\claim{enclave-isolation}. \\
    A9 & Close time tracks wall-clock time; ledgers close on average no
          faster than the target close time, from which the wallet sizes
          Perch's 7-day delay floor\claim{ledger-timing}. \\
    A10 & Mail to the recovery address reaches only the enclave: the
          operator does not change the recovery domain's MX, MTA-STS
          policy~\cite{rfc8461} or the mail exchanger's TLS
          keys\claim{recovery-domain-routing}. \\
    A11 & Users run the published wallet build\claim{published-wallet-build}. \\
    A12 & Under email plus delay, the owner learns of an authorized attempt
          within the delay\claim{alert-liveness}. \\
    \bottomrule
  \end{tabularx}
\end{table}

\paragraph{Goals.} The email factor should satisfy four properties, stated
formally in \S\ref{sec:properties}.
\begin{description}
  \item[G1 Soundness.] An accepted approval implies a DKIM signature by a
    registered key for the enrolled domain over an email whose sender is the
    committed address and whose content binds the exact recovery statement:
    its account, new key, nonce and expiry.
  \item[G2 Privacy.] The chain's state, the transactions and the proofs
    reveal nothing about the address beyond the salted commitment.
  \item[G3 Replay and front-running resistance.] An approval is bound to one
    account, one statement and one nonce, and a copy submitted by anyone
    records the same single vote.
  \item[G4 One-factor safety.] No sequence of email approvals alone
    completes a recovery: completion still needs the delay and the policy's
    quorum.
\end{description}
G4 is the precise sense of ``email is never sufficient alone'': it never
skips the delay. Under the lost-key profile with the email guardian as the
only one needed (\emph{email plus delay}: at least seven days, and a second
alert channel registered when the wallet builds the configuration), email
alone authorizes, and what stops a wrongful recovery is the owner's
cancellation within the delay, which rests on A12. The wallet checks the
channel when it builds a configuration; if the channel later lapses or goes
unacknowledged, the chain still accepts the email vote alone, and the wallet
tells the owner and offers to add a device or another helper.

Off chain the goal is scoped. The wallet tells users exactly
\PrivacyClaim, a claim against the operator's servers and staff that holds
under A8, A10 and A11; it is not a claim that the address can never
leak\claim{scoped-privacy-claim}. The user's mail provider and the
alert-sending service necessarily see the address. The provider sees the
blinded Subject, whose content names no attempt; the timing of the email
and of the vote that follows it can (\S\ref{sec:limits}).

\paragraph{Beyond the protocol.} G1 binds the email to the statement, not to
the sender's informed consent. The social engineer needs no fake page: the
real recovery page on its own device opens a lost-key attempt for its key
(opening is permissionless) and shows a Subject the user cannot check, and
since alert codes derive from $\code$, the attacker holding it sees the code
the alerts carry. When the victim acts, typing the dictated Subject into
their own mail client, no text of the operator's reaches them: the prepared
body's warning and the warning where the code is typed render on the
attacker's screen and aid honest users only. The alerts reach the victim
after the vote, and each required alert says that whoever told the reader to
expect it, or to ignore it, is trying to take the
account\claim{phishing-defences}. What remains is the delay with the owner's
cancellation, the push to the device set up for alerts, the wallet's own read
of open attempts, and other guardians; a \texttt{Re:} reply to a phishing
message is unprovable. Whether the recovery address's mail exchanger should
reply to each approval's sender with a warning, with \texttt{To:} bound
where signed so the reply cannot be skipped, is open
(Appendix~\ref{app:open}). Two risks remain under email plus
delay: a phished user who has lost every device cannot cancel, and with
alerts going to a stolen inbox and no push device reachable, only opening the
wallet shows the attempt (\S\ref{sec:eval-attacks}). Whether users act on
what reaches them is not measured\claim{user-behaviour}.

\paragraph{Which losses the factor covers.} Platforms that sync passkeys
sync the wallet's by default, so losing every device loses such a passkey
only with the platform account or its keychain's unlock
secret\claim{passkey-sync}. For a Gmail or iCloud user the platform account
can also hold the factor's inbox (Table~\ref{tab:losses}); enrollment notes
such an inbox and recommends another inbox or
helper\claim{enrollment-safeguards}.

\begin{table}[t]
  \caption{When the email factor adds recovery coverage, by passkey and
    loss, assuming the inbox's provider does not require a lost device to
    sign in.}
  \label{tab:losses}
  \small
  \begin{tabularx}{\columnwidth}{@{}>{\raggedright\arraybackslash}p{0.2\columnwidth}>{\raggedright\arraybackslash}X>{\raggedright\arraybackslash}X>{\raggedright\arraybackslash}X@{}}
    \toprule
    Lost & Device-bound passkey & Synced, inbox at the sync provider & Synced, inbox elsewhere \\
    \midrule
    Every device & covers & not needed: sync restores the passkey & not needed \\
    Platform account & --- & adds nothing: the inbox is lost too & covers \\
    Keychain unlock secret only & --- & covers & covers \\
    \bottomrule
  \end{tabularx}
\end{table}

\paragraph{Non-goals.} The design does not hide that an account has an email
guardian, that a recovery happened, or, from the mail provider, that its user
sent a recovery approval and when. An adversary with the inbox and the
account code who meets the quorum with other factors, or outlasts the delay
unseen, recovers the account: that is the policy's semantics. A compromised
user device, custom mail domains and a provider that signs mail its user did
not send are out of scope.
